\label{chap.2}
\section{Attacks and Benchmarks}
Greshake et al.~\cite{greshake2023} were the first to describe indirect prompt injection. They showed that content retrieved by an LLM application can be used to steal data, spread injections by itself and control the application from outside. Liu et al.~\cite{liu2023} later built a black-box attack framework and broke into 31 of 36 real LLM-integrated applications.

Several benchmarks came out after that. InjecAgent~\cite{zhan2024} has 1,054 test cases with 17 user tools and 62 attacker tools, split into direct harm and data stealing. AgentDojo~\cite{debenedetti2024} is a dynamic environment with 97 realistic tasks and 629 security cases for banking, email, travel and workspace agents. BIPIA~\cite{yi2023} covers indirect injection over several application tasks, and the Agent Security Bench (ASB)~\cite{zhang2024asb} tests agents that have many tools and a memory. In most of these benchmarks the injected instruction is plain text inside a tool output. So they check whether a defence can spot a command, but not whether it can follow an item from one hop to the next.

\section{Defences}
Task Shield~\cite{jia2024} treats security as a question of task alignment and checks whether each instruction or tool call actually serves the user's original goal. AgentWatcher~\cite{wang2026} links an action back to the parts of the context that caused it and lets a monitor LLM reason over explicit rules. LlamaFirewall~\cite{llamafirewall2025} combines a prompt-injection classifier with AlignmentCheck, where an LLM reviews the agent's whole trajectory against the user's goal. Li and Zhao~\cite{li2026} studied agents trained with such defences and found that they complete far fewer tasks. Geng et al.~\cite{geng2026} reviewed 128 studies and pointed out that common evaluation setups are limited and that new or adaptive attacks are rarely tested.

\section{Multi-Agent Propagation}
Prompt Infection~\cite{lee2024} shows injections that copy themselves from agent to agent. The G-Safeguard system~\cite{gsafeguard2025} builds a graph of the messages between agents and uses graph neural networks to find suspicious nodes and cut edges. Both make it clear that spreading across agents is a real threat. Neither offers a light monitor that sits inline, keeps provenance for each session, and is tested on fragment attacks with matching harmless twins.

\section{General-Purpose Guardrails}
Llama Guard~\cite{inan2023} decides whether a conversation turn is safe. NeMo Guardrails~\cite{rebedea2023} lets developers add programmable rails, one of which asks an LLM to check the input. Prompt Guard 2~\cite{promptguard2025} and the ProtectAI DeBERTa classifier~\cite{protectai2024} look for injection-like text in a single input. All of them look at one message at a time, and Chapter~6 shows how much that costs on multi-hop attacks.

\section{Comparison}
\begin{table}[H]
\centering
\small
\begin{tabular}{p{3.2cm}p{4.2cm}p{5.6cm}}
\toprule
Work & Approach & Limit addressed by MAPIS\\
\midrule
Task Shield & Per-action task alignment by an LLM & Relies on the backbone LLM's judgement; evaluated on one benchmark\\
AgentWatcher & Attribution plus monitor LLM & Adds LLM inference to every check\\
LlamaFirewall & Classifier plus trajectory audit by an LLM & Heavy and post-hoc: an LLM call over the whole trajectory\\
G-Safeguard & GNN over the agent graph & Not a lightweight inline stateful tap\\
Llama Guard, Prompt Guard 2, ProtectAI & Single-message classifiers & No session memory\\
\bottomrule
\end{tabular}
\caption{Related defences and the limit that MAPIS addresses.}
\end{table}

\section{Research Gap}
The gap we address is a monitor that is light and real-time, that (i) keeps session state across agent hops, (ii) is tested on fragment attacks together with harmless twins that look the same message by message, and (iii) works with more than one agent framework. MAPIS is not the first work on injection in multi-agent systems. What is new is putting these three things together at a tap that does not depend on the framework.
